\section{Related directions and the open target}
Connections between resource-bounded descriptions and lower bounds
must retain their models, time scales, promises, and hypotheses.

Chapman and Williams \cite{ChapmanWilliams2015} characterize
\(\NP\nsubseteq\Ppoly\) through natural properties designed against
SAT-solving circuits with specified checking behavior.
This concerns a specialized class of solvers, not arbitrary
properties separating SAT from arbitrary circuits. It is a
literature direction, not a lower-bound method supplied here.

Oliveira, Pich, and Santhanam \cite{OliveiraPichSanthanam2021},
Theorem~1.4, prove the following hardness-magnification implication.
There is a universal \(c\geq1\) such that, if some \(\varepsilon>0\)
satisfies, for every sufficiently small constant \(\beta>0\),
\[
 \mathsf{Gap\text{-}MCSP}
 \left[\frac{N^\beta}{c\log_2 N},N^\beta\right]
 \notin\SIZE[N^{1+\varepsilon}],
\]
then \(\NP\nsubseteq\Ppoly\).
Here \(N\) is the length of the entire input truth table. The
promise distinguishes circuit complexity at most the first threshold
from complexity at least the second.
A restricted diagram bound for \(T_n\) does not satisfy this
antecedent. An application requires a separate resource-preserving
reduction.

Kabanets and Kolokolova \cite{KabanetsKolokolova2026} characterize
specified chain rules for time-bounded Kolmogorov complexity.
The March 13, 2026 revision states a direct equivalence between a
conditional chain rule with sublinear error and \(\PP=\NP\),
\emph{assuming} NP-hardness of the stated
\(\mathsf{GapMcK}^{t}\mathsf P\) problem.
It also gives unconditional characterizations involving other
complexity assumptions. We neither remove the extra hypothesis nor
identify their full-output conditional measure with local access.

\begin{conjecture}[Research target \claimid{Q01}]
\label{conj:main}
For every \(a,b\in\mathbb N_{>0}\), there are infinitely many \(n\)
such that \(\loc_{n^b}(T_n)>n^a\).
\end{conjecture}
By Corollary~\ref{cor:target}, this is the established open target
\(\SAT\notin\Ppoly\) in local-description language, not a new
conjecture or a proof advance.

An intermediate result must name its restricted model and identify
the changes of representation it tolerates. Candidate directions
include order-independent residual bounds for specified diagram
classes, parameter-checked reductions to a meta-complexity promise
problem, and proof-assistant formalization of the equivalence.
These tasks remain open in this project.
A bound \(|d|t^\alpha\geq n^c\) for fixed \(c\), on its own, still
permits polynomial description length and polynomial time.
